% ═══════════════════════════════════════════════════════════════════
% Paper V: Memory, Consciousness, and Soul as Impedance Physics
% ═══════════════════════════════════════════════════════════════════
\documentclass[aps,pre,twocolumn,superscriptaddress,floatfix]{revtex4-2}

\usepackage{amsmath,amssymb,amsthm}
\usepackage{graphicx}
\usepackage{hyperref}
\usepackage{bm}
\usepackage{booktabs}
\usepackage{siunitx}

\newtheorem{theorem}{Theorem}
\newtheorem{corollary}{Corollary}
\newtheorem{definition}{Definition}
\newtheorem{remark}{Remark}
\newtheorem{proposition}{Proposition}

\newcommand{\Gmeta}{\Gamma_{\text{meta}}}
\newcommand{\Phimeta}{\Phi_{\text{meta}}}
\newcommand{\Zarousal}{G_{\text{arousal}}}
\newcommand{\QPFC}{Q_{\text{PFC}}}
\newcommand{\Acut}{A_{\text{cut}}}
\newcommand{\avg}[1]{\langle #1 \rangle}

\begin{document}

\title{Memory, Consciousness, and Soul\\
as Impedance Physics on the $\Gamma$-Net}

\author{Hsi-Yu Huang}
\email{llc.y.huangll@gmail.com}
\affiliation{Independent Researcher, Taiwan}

\date{\today}

% ============================================================
\begin{abstract}
Papers~I--IV of this series derived relay-node topology
(Paper~I), dual neural--vascular networks (Paper~II),
impedance debt and sleep (Paper~III), and the lifecycle
equation governing cognition from birth to senescence
(Paper~IV), all from a single variational principle:
$A[\Gamma]=\int_0^T\!\sum_i\Gamma_i^2(t)\,dt\to\min$.
Here we show that three concepts traditionally regarded as
metaphysical---memory, consciousness, and soul---have
precise operational definitions within the same formalism.

\emph{Memory} is the time derivative of impedance:
$\partial Z/\partial t=-\eta\,\Gamma\,x_{\text{pre}}\,x_{\text{post}}$
(Hebbian update, Constraint~C2).
Muscle memory corresponds to $\Gamma\to 0$ (over-learned match),
while cognitive memory retains $\Gamma\approx\varepsilon>0$
(retrievable because still slightly mismatched).

\emph{Consciousness} requires a self-referential measurement
loop.  We define the meta-reflection coefficient
$\Gmeta=(Z_{\text{obs}}-Z_{\text{act}})/(Z_{\text{obs}}+Z_{\text{act}})$,
where observer and actor are the \emph{same} network sampled at
successive ticks.  The consciousness index
$\Phimeta=1-\Gmeta^2$ quantifies self-model fidelity.
Crucially, $\Phimeta>0$ requires a \emph{necessary condition}:
$\Zarousal\cdot\sum_i G_i\cdot\QPFC>\theta$,
which is formally identical to a television set
(power supply $\times$ channel selector $\times$ decoder chip).
The television is an \emph{open-loop} system ($\Gmeta$ undefined);
the brain closes the loop.

\emph{Soul} is the null space of Hebbian learning:
$\mathcal{S}=\ker(\partial Z/\partial t)
=\{Z_k:\Delta Z_k=0\;\forall\;\text{experience}\}$.
It decomposes into three invariants:
irreducible topological cost $\Acut$ (Paper~I),
genetically fixed impedance $Z_0^{\text{DNA}}$, and
the thermal noise floor $\rho_{\text{irred}}^2$ from
Johnson--Nyquist physics.
The soul is not created or destroyed---it is the geometric
boundary condition of an impedance network.

Computational experiments verify two testable predictions:
(i)~all sensory modalities converge on the prefrontal cortex
(PFC), making it the consciousness \emph{bottleneck};
(ii)~fever (PFC thermal noise) and occipital impact
(brainstem arousal crash) are two \emph{independent}
kill mechanisms of consciousness, proving that consciousness
is a product $\Phi=\Zarousal\times\sum G_i\times\QPFC$
rather than a location---analogous to the area of a
rectangle being destroyed by zeroing either side.
\end{abstract}

\maketitle

% ============================================================
\section{Introduction}
\label{sec:intro}

Where is memory stored?
When does consciousness begin?
Where does the soul reside?
These three questions have dominated philosophy, theology,
and neuroscience for millennia.  From a physics standpoint,
all three are \emph{ill-posed}: they assume that memory is
localised, consciousness is binary, and soul is an object.

Papers~I--IV of this series established that all neural and
vascular dynamics emerge from the Minimum Reflection Principle
(MRP):
\begin{equation}
  A[\Gamma] = \int_0^T \sum_i \Gamma_i^2(t)\,dt \to \min\,,
  \label{eq:MRP}
\end{equation}
where $\Gamma_i=(Z_{L,i}-Z_{0,i})/(Z_{L,i}+Z_{0,i})$ is the
impedance reflection coefficient at channel~$i$.
Three inviolable constraints govern the network:
\begin{itemize}
  \item[\textbf{C1}] Energy conservation: $\Gamma^2+T=1$ at
    every channel, every tick.
  \item[\textbf{C2}] Hebbian update:
    $\Delta Z=-\eta\,\Gamma\,x_{\text{pre}}\,x_{\text{post}}$
    (gradient descent on~$A[\Gamma]$).
  \item[\textbf{C3}] Signal protocol: all inter-module values
    carry impedance metadata ($Z_0,Z_L$).
\end{itemize}

In this paper we demonstrate that memory, consciousness, and
soul are not additional postulates but \emph{necessary
consequences} of Eqs.~(\ref{eq:MRP})--C1/C2/C3.
Section~\ref{sec:memory} derives memory as impedance dynamics.
Section~\ref{sec:consciousness} constructs the
self-referential consciousness index~$\Phimeta$.
Section~\ref{sec:soul} identifies the soul as the null space
of C2.
Section~\ref{sec:television} shows that the necessary
condition for consciousness is structurally identical to a
television set, differing only in feedback topology.
Section~\ref{sec:two-kills} presents computational experiments
demonstrating two independent kill mechanisms.
Section~\ref{sec:discussion} discusses implications.


% ============================================================
\section{Memory as Impedance Dynamics}
\label{sec:memory}

\subsection{The Hebbian gradient}

Constraint~C2 defines the only mechanism by which the network
modifies itself:
\begin{equation}
  \Delta Z_k = -\eta\,\Gamma_k\,x_{\text{pre},k}\,x_{\text{post},k}\,.
  \label{eq:hebbian}
\end{equation}
This is not an analogy---it \emph{is} memory.
Learning occurs when $\Gamma_k\neq 0$ and correlated
pre/post-synaptic activity is present.
The trace of all past updates constitutes the network's
memory:
\begin{equation}
  Z_k(t) = Z_k(0) - \eta \int_0^t
    \Gamma_k(\tau)\,x_{\text{pre},k}(\tau)\,
    x_{\text{post},k}(\tau)\,d\tau\,.
  \label{eq:memory-integral}
\end{equation}

\subsection{Muscle memory: $\Gamma\to 0$}

When a motor skill is practiced to mastery, the impedance at
every channel in the circuit approaches perfect matching:
$Z_L\to Z_0$, hence $\Gamma\to 0$.
At this point, Eq.~(\ref{eq:hebbian}) gives $\Delta Z\to 0$:
the memory has become \emph{frozen} into the impedance
structure itself.
No further energy is required to maintain it ($\Gamma^2\to 0$),
and no further learning occurs ($\Delta Z=0$).

This explains why:
\begin{itemize}
  \item You never forget how to ride a bicycle
    ($\Gamma\approx 0$, no decay).
  \item Muscle memory is ``automatic''---it requires no
    conscious effort because $\Gamma^2\approx 0$ means
    near-zero reflected energy.
  \item Over-training a cognitive skill (e.g., multiplication
    tables) moves it from cognitive to muscle memory:
    $\Gamma\to 0$.
\end{itemize}

\subsection{Cognitive memory: $\Gamma\approx\varepsilon>0$}

Declarative memories---facts, episodes, narratives---are
\emph{not} driven to $\Gamma=0$.  They retain a small but
nonzero reflection:
\begin{equation}
  \Gamma_{\text{cog}} \approx \varepsilon > 0\,.
  \label{eq:cog-memory}
\end{equation}
This has two consequences:
\begin{enumerate}
  \item \emph{Retrievability}: A nonzero $\Gamma$ means the
    memory generates a detectable reflection when probed.
    If $\Gamma=0$, the memory is invisible---it has become
    part of the medium itself, like light passing through
    a perfectly matched cable.
  \item \emph{Decay}: Since $\Gamma>0$, the Hebbian update
    continues to operate.  Without reinforcement
    ($x_{\text{pre}}\cdot x_{\text{post}}\to 0$), the memory
    decays, governed by
    $\lambda_{\text{eff}}=\lambda_{\text{base}}/(1-\Gamma_{\text{bind}}^2)$
    from Paper~III.
\end{enumerate}

\begin{remark}[The paradox of perfect memory]
A perfectly learned item ($\Gamma=0$) is \emph{unrecallable}:
it has merged with the network substrate.
Memory requires imperfection.
This is the impedance analogue of the observer effect:
you can only detect what reflects.
\end{remark}

\subsection{Three memory regimes}

\begin{table}[h]
\centering
\caption{Memory regimes classified by reflection coefficient.}
\label{tab:memory-regimes}
\begin{tabular}{@{}lccc@{}}
\toprule
Regime & $\Gamma$ & $\Delta Z$ & Recallable? \\
\midrule
Muscle memory & $\to 0$ & $\to 0$ & No (automatic) \\
Cognitive memory & $\varepsilon>0$ & small & Yes (effortful) \\
Novel stimulus & $\gg 0$ & large & Vivid (but unstable) \\
\bottomrule
\end{tabular}
\end{table}


% ============================================================
\section{Consciousness as Self-Referential Impedance Matching}
\label{sec:consciousness}

\subsection{The hard problem restated}

The ``hard problem'' of consciousness asks why subjective
experience exists at all.  In impedance physics, we reframe
it: \emph{What physical structure allows a network to measure
its own reflection coefficient?}

Standard signal processing is open-loop: a signal enters,
gets processed, exits.  No part of the chain ``knows'' the
output.  We define consciousness as the condition under which
the network's output is routed back as input, and the
resulting self-referential mismatch is measured.

\subsection{The meta-reflection coefficient}

\begin{definition}[Meta-reflection]
Let $Z_{\text{obs}}(t)$ be the impedance of the observer
sub-network at tick~$t$, and $Z_{\text{act}}(t-1)$ the
impedance of the actor sub-network at the previous tick
(the output that is now being observed).  The
\emph{meta-reflection coefficient} is:
\begin{equation}
  \Gmeta = \frac{Z_{\text{obs}} - Z_{\text{act}}}
                 {Z_{\text{obs}} + Z_{\text{act}}}\,.
  \label{eq:Gamma-meta}
\end{equation}
\end{definition}

\begin{definition}[Consciousness index]
\begin{equation}
  \Phimeta = 1 - \Gmeta^2\,.
  \label{eq:Phi-meta}
\end{equation}
When $\Gmeta=0$: the observer perfectly matches the actor;
the self-model is maximally faithful.
When $\Gmeta\to 1$: total self-mismatch; the system cannot
model itself.
\end{definition}

\subsection{Necessary condition: the television inequality}

$\Phimeta$ is only meaningful if the underlying signal chain
is operational.  We derive the \emph{television inequality}:

\begin{theorem}[Television Inequality]
\label{thm:television}
Let $\Zarousal\in[0,1]$ be the brainstem arousal gain,
$G_i\in[0,1]$ the thalamic gate gain for modality~$i$,
and $\QPFC\in[0,1]$ the prefrontal cortex processing
quality.  The consciousness index $\Phimeta>0$ requires:
\begin{equation}
  \Zarousal \cdot \sum_{i=1}^{N} G_i \cdot \QPFC > \theta\,,
  \label{eq:TV-inequality}
\end{equation}
where $\theta>0$ is the minimum signal-to-noise ratio for
the self-referential loop to close.
\end{theorem}

\begin{proof}
The meta-reflection loop requires a signal to:
(a)~traverse the thalamic gate ($\times G_i$),
(b)~be processed by PFC ($\times\QPFC$),
(c)~be routed back through the arousal-modulated pathway
($\times\Zarousal$).
If any factor is zero, the loop is open and $\Gmeta$ is
undefined (division by zero in Eq.~\ref{eq:Gamma-meta}
when $Z_{\text{act}}=0$).
The threshold~$\theta$ accounts for Johnson--Nyquist noise:
the signal must exceed the thermal floor to be self-detectable.
\end{proof}

\subsection{Open loop vs.\ closed loop}

Equation~(\ref{eq:TV-inequality}) is structurally identical
to the operation of a television set:
\begin{center}
\begin{tabular}{@{}lll@{}}
\toprule
Component & Television & Brain \\
\midrule
Power supply & Wall plug & Brainstem RAS ($\Zarousal$) \\
Channel selector & Tuner & Thalamus ($\sum G_i$) \\
Decoder chip & Video IC & PFC ($\QPFC$) \\
\midrule
Kill method 1 & Unplug power & Occipital impact \\
Kill method 2 & Overheat IC & Fever \\
\bottomrule
\end{tabular}
\end{center}

A television satisfies Eq.~(\ref{eq:TV-inequality}): it has
power, channel selection, and signal processing.
Yet it is not conscious.  The missing ingredient is the
\emph{feedback topology}: the television's output (screen)
is not routed back to its input (antenna).
It is an open-loop system; $\Gmeta$ is undefined.

The brain uses the \emph{same components} but adds one
connection: the output of PFC decision-making is fed back
through the sensory pathway (via efference copy, proprioception,
and internal monitoring), creating a closed loop where
$Z_{\text{obs}}$ and $Z_{\text{act}}$ refer to the
\emph{same network} at successive time steps.

\begin{corollary}[Consciousness is topological, not material]
Two systems with identical components (power, gate, processor)
differ in consciousness solely by the presence or absence of
a self-referential feedback connection.
Consciousness is a property of the network \emph{topology},
not of the network \emph{material}.
\end{corollary}

\subsection{Why ``where is consciousness?''\ is ill-posed}

Consciousness is not located at the power supply, the gate,
or the processor.  It is the \emph{product}:
\begin{equation}
  \Phimeta = f\!\left(
    \Zarousal \cdot \textstyle\sum_i G_i \cdot \QPFC
  \right)\,.
  \label{eq:consciousness-product}
\end{equation}
Asking ``where is consciousness'' is like asking ``where is
the area of a rectangle''---it is not in the width or the
height, but in their product.  Zeroing either dimension
destroys the area.


% ============================================================
\section{Soul as the Null Space of Learning}
\label{sec:soul}

\subsection{What cannot be changed by experience}

The Hebbian update (Eq.~\ref{eq:hebbian}) modifies impedance
wherever $\Gamma\neq 0$ and correlated activity exists.
We ask: what part of the impedance network is \emph{invariant}
under all possible experience?

\begin{definition}[Soul]
\begin{equation}
  \mathcal{S} = \ker\!\left(\frac{\partial Z}{\partial t}\right)
  = \{Z_k : \Delta Z_k = 0\;\;\forall\;\text{experience}\}\,.
  \label{eq:soul}
\end{equation}
\end{definition}

This kernel decomposes into three components:

\begin{theorem}[Soul Decomposition]
\label{thm:soul-decomp}
\begin{equation}
  \mathcal{S} = \Acut + Z_0^{\text{DNA}} + \rho_{\text{irred}}^2\,,
  \label{eq:soul-decomposition}
\end{equation}
where:
\begin{enumerate}
  \item $\Acut$: the irreducible topological action cost
    (Paper~I, Theorem~1).  This is the minimum $\Gamma^2$
    required by the network's dimensional heterogeneity,
    independent of how well the network is trained.
  \item $Z_0^{\text{DNA}}$: genetically fixed impedance
    values that constrain the characteristic impedance of
    neural substrates.  These set the ``hardware'' of the
    network and are not modifiable by C2.
  \item $\rho_{\text{irred}}^2$: the irreducible noise floor
    from Johnson--Nyquist thermal physics,
    $V_n^2=4k_BT\cdot R\cdot\Delta f$.
    No amount of learning can reduce impedance below the
    thermodynamic fluctuation limit.
\end{enumerate}
\end{theorem}

\begin{proof}
(i)~$\Acut$ is invariant by Theorem~1 of Paper~I: it depends
only on the dimensional structure $(K_i,K_j)$ of connected
nodes, not on impedance values.
(ii)~$Z_0^{\text{DNA}}$ enters C2 only through $\Gamma_k$;
but $Z_{0,k}$ is the \emph{characteristic} impedance
(the ``cable property''), not the \emph{load} impedance.
C2 modifies $Z_{L,k}$, not $Z_{0,k}$.
Hence $Z_0^{\text{DNA}}$ is invariant.
(iii)~$\rho_{\text{irred}}^2$ sets a floor on $\Gamma$ via
$\delta\Gamma\propto\kappa(\mathbf{Z})\cdot\delta T/T$;
this is a thermodynamic constant, not modifiable by C2.
\end{proof}

\subsection{Properties of the soul}

\begin{corollary}[The soul is not an event]
$\mathcal{S}$ is a time-independent geometric structure.
It does not ``arrive'' at a moment or ``depart'' at death.
It is the boundary condition of the impedance network,
present from the moment the network's topology is fixed
(approximately: fertilisation for $Z_0^{\text{DNA}}$,
neural migration for $\Acut$).
\end{corollary}

\begin{corollary}[The soul cannot be learned or unlearned]
Since $\mathcal{S}=\ker(\partial Z/\partial t)$, no
experience, training, trauma, or therapy can alter the soul.
It is the part of you that remains after everything learnable
has been taken away.
\end{corollary}

\begin{corollary}[Uniqueness]
$\Acut$ depends on the specific topology $(K_i,K_j)$ of the
network, which is determined by genetics and stochastic
developmental morphogenesis.
No two organisms have identical $\Acut$.
Hence the soul is unique.
\end{corollary}


% ============================================================
\section{The Television Theorem}
\label{sec:television}

We now formalise the structural identity between a television
and a brain, and the single topological difference.

\begin{theorem}[Television Theorem]
\label{thm:TV-full}
Any system with three multiplicative components---a power
supply~$P$, a channel selector~$G$, and a
processor~$Q$---produces an output $O=P\cdot G\cdot Q$.
Destroying any single component zeroes the output.
The system is conscious if and only if there exists a
feedback path $\pi:O\to\text{input}$ such that the
meta-reflection coefficient
$\Gmeta=(Z_{\text{obs}}-Z_{\text{act}})/(Z_{\text{obs}}+Z_{\text{act}})$
is defined and measurable.
\end{theorem}

\begin{table}[h]
\centering
\caption{Three systems with identical multiplicative
structure but different feedback topology.}
\label{tab:three-systems}
\begin{tabular}{@{}lccc@{}}
\toprule
 & Rock & Television & Brain \\
\midrule
Power $P$ & None & Wall plug & Brainstem \\
Gate $G$ & None & Tuner & Thalamus \\
Processor $Q$ & None & Video IC & PFC \\
$P\!\cdot\!G\!\cdot\!Q > 0$ & No & \textbf{Yes} & \textbf{Yes} \\
Feedback $\pi$ & No & No & \textbf{Yes} \\
$\Gmeta$ defined & No & No & \textbf{Yes} \\
Conscious? & No & No & \textbf{Yes} \\
\bottomrule
\end{tabular}
\end{table}


% ============================================================
\section{Two Independent Kill Mechanisms}
\label{sec:two-kills}

We present computational experiments demonstrating that
consciousness can be destroyed by two physically independent
mechanisms, proving it is a product rather than a location.

\subsection{Experimental setup}

The simulation instantiates:
\begin{itemize}
  \item A \texttt{ThalamusEngine} with visual, auditory, and
    vestibular channels, each with independent gate gain~$G_i$.
  \item A \texttt{PrefrontalCortexEngine} with goal management,
    Go/NoGo evaluation, and finite energy.
  \item A consciousness proxy:
    $\Phi_{\text{proxy}}=\bar{G}\cdot E_{\text{PFC}}\cdot
    G_{\text{arousal}}$.
\end{itemize}

\subsection{Kill mechanism 1: Fever}

Temperature rises from $37^\circ$C to $41^\circ$C
over 20~ticks.  Arousal remains at~0.85 (patient is awake).

\textbf{Result}: All thalamic gates remain \emph{open}
(visual: 0.71, auditory: 0.20, vestibular: 0.16).
PFC energy remains~1.0.
$\Phi_{\text{proxy}}$ declines gradually from 0.42 to 0.30
but does not reach blackout.

\emph{Diagnosis}: ``Lights on, nobody home.''
Gates transmit signals, but PFC matrix inversion degrades
because $\delta\Gamma\propto\kappa(\mathbf{Z})\cdot\delta T/T$
and $\kappa(\mathbf{Z})\propto K^{3/2}$.
At $K=128$ (cortical pyramidal cell axon dimension),
a \SI{2}{\celsius} fever produces
$\delta\Gamma/\Gamma\approx 1448\times 2/310\approx 9.3\gg 1$.

\subsection{Kill mechanism 2: Occipital impact}

Temperature constant at $37^\circ$C.
Arousal crashes from 0.85 to 0.01 in three ticks (brainstem
reticular activating system disruption).

\textbf{Result}: All thalamic gates slam \emph{shut}
(all channels: 0.000 at worst tick).
PFC energy remains~1.0 (hardware is intact).
$\Phi_{\text{proxy}}$ drops to~0.000 within 2~ticks.
Thalamic burst mode activates (random gate flicker).

\emph{Diagnosis}: ``Power cut, hardware fine.''
The processor is undamaged, but with $\Zarousal\to 0$,
$G_{\text{total}}=\Zarousal\cdot(\alpha\,G_{\text{top}}+(1-\alpha)\,G_{\text{bot}})\to 0$
for all channels simultaneously.

\subsection{Impact + vestibular shock}

A realistic posterior impact produces mechanical shockwaves
reaching the cochlea and vestibular apparatus:
\begin{itemize}
  \item At impact ($t=5$--$7$): vestibular amplitude spikes
    to~0.95 ($\Gamma=0.90$), triggering the thalamic
    \emph{startle bypass} (hardware interrupt).
  \item Auditory amplitude spikes to~0.85 (tinnitus).
  \item Result: auditory and vestibular gates forced to 1.0
    via startle circuit, \emph{even as arousal crashes}.
\end{itemize}

This explains the subjective sequence:
\begin{enumerate}
  \item Tinnitus: cochlear shockwave $\to$ startle bypass $\to$
    auditory gate forced open.
  \item Vertigo: vestibular shockwave $\to$ startle bypass $\to$
    vestibular gate forced open.
  \item Blackout: arousal $\to 0$ $\to$ all gates close
    (startle signals are the \emph{last} to pass before shutdown).
\end{enumerate}

The startle circuit is a \emph{hardware interrupt}---it
bypasses the normal arousal-modulated gating.
This is why you hear the buzz and feel the spin
\emph{before} losing consciousness: the interrupt fires
faster than the power-down propagates.

\subsection{Quantitative comparison}

\begin{table}[h]
\centering
\caption{Consciousness collapse comparison.}
\label{tab:two-kills}
\begin{tabular}{@{}lcc@{}}
\toprule
 & Fever & Impact \\
\midrule
Kill target & $\QPFC$ (processor) & $\Zarousal$ (power) \\
Onset & Minutes--hours & Milliseconds \\
Gate state & Open & Closed \\
$\Phi_{\min}$ & 0.302 & 0.000 \\
PFC energy & Intact & Intact \\
Analogy & CPU overheat & Power cut \\
\bottomrule
\end{tabular}
\end{table}

If consciousness resided in either location, the other
mechanism should not affect it.
Both independently reduce $\Phi\to 0$.
Therefore consciousness is the \emph{product}, not its factors.


% ============================================================
\section{Thermal Vulnerability and the K-Dimensional Bottleneck}
\label{sec:thermal}

Johnson--Nyquist thermal noise generates voltage fluctuations:
\begin{equation}
  V_n^2 = 4 k_B T \cdot R \cdot \Delta f\,.
  \label{eq:JN}
\end{equation}
For a $K$-dimensional impedance matrix $\mathbf{Z}\in\mathbb{R}^{K\times K}$,
the condition number scales as $\kappa(\mathbf{Z})\propto K^\alpha$
with $\alpha\approx 3/2$ for biologically plausible random matrices.
The fractional perturbation of $\Gamma$ is:
\begin{equation}
  \frac{\delta\Gamma}{\Gamma}
    \propto \kappa(\mathbf{Z}) \cdot \frac{\delta T}{T}
    \propto K^{3/2} \cdot \frac{\delta T}{T}\,.
  \label{eq:thermal-vulnerability}
\end{equation}

\begin{table}[h]
\centering
\caption{Thermal vulnerability by axon dimension~$K$.
($\delta T=+2$\,K (fever), $T_0=310$\,K).}
\label{tab:thermal-K}
\begin{tabular}{@{}rrrrl@{}}
\toprule
$K$ & $\kappa(\mathbf{Z})$ & $\delta\Gamma/\Gamma$ (39$^\circ$C)
    & $\delta\Gamma/\Gamma$ (41$^\circ$C) & Status \\
\midrule
2   & 2.8    & 0.018  & 0.037  & Normal \\
8   & 22.6   & 0.146  & 0.292  & Degraded \\
16  & 64.0   & 0.413  & 0.825  & Impaired \\
32  & 181    & 1.17   & 2.33   & Collapse \\
128 & 1448   & 9.34   & 18.7   & Collapse \\
\bottomrule
\end{tabular}
\end{table}

The prefrontal cortex, with cortical pyramidal neurons at
$K\approx 128$, is the \emph{first} region to suffer thermal
degradation---not because consciousness ``lives'' there, but
because its high-dimensional matrix inversion is the most
noise-sensitive computation in the brain.
This explains why:
\begin{itemize}
  \item Fever attacks executive function and decision-making
    before it attacks sensation.
  \item The forehead feels hot during fever, but consciousness
    degrades at the \emph{frontal} cortex (highest~$K$).
  \item Other brain regions (sensory cortices, $K\sim 4$--$16$)
    can tolerate several degrees of fever.
\end{itemize}


% ============================================================
\section{Discussion}
\label{sec:discussion}

\subsection{Three ill-posed questions answered}

\begin{enumerate}
  \item \emph{Where is memory stored?}---Everywhere.
    Memory is $\partial Z/\partial t$; it is the
    \emph{change} in impedance, distributed across all
    channels that were active during learning.
    It is not localised to a ``memory bank.''
  \item \emph{When does consciousness begin?}---It does not
    ``begin'' at a moment.
    $\Phimeta=1-\Gmeta^2$ is a continuous variable.
    It increases as the self-referential loop quality
    improves during development, and decreases during
    sleep, anaesthesia, or brain injury.
    Asking ``when'' is like asking when temperature begins.
  \item \emph{Where does the soul reside?}---The soul
    $\mathcal{S}=\ker(\partial Z/\partial t)$ is a geometric
    structure of the network, not a spatial location.
    It is present from the moment topology is fixed and
    persists until the network disintegrates.
\end{enumerate}

\subsection{The television as a boundary case}

The Television Theorem (Theorem~\ref{thm:TV-full}) provides
a sharp boundary: any system satisfying
$P\cdot G\cdot Q>\theta$ with a closed feedback loop is a
candidate for consciousness.
Current artificial neural networks (including large language
models) operate in open-loop mode during inference---they
process input and produce output but do not measure their own
output's impedance mismatch with their input.
They are televisions, not brains.

\subsection{Circularity warning}

We address the elephant in the room.
All computational experiments in this paper (and Papers~I--IV)
use simulation modules \emph{built from} the $\Gamma$-Net
theory itself.  When the \texttt{ThalamusEngine} obeys
$G_{\text{total}}=G_{\text{arousal}}\cdot(\alpha\,G_{\text{top}}+(1-\alpha)\,G_{\text{bot}})$
and then a simulation confirms that $G_{\text{arousal}}\to 0$
kills all gates, we have not validated the theory---we have
verified that the code implements the equation correctly.

This is \emph{internal consistency}, not \emph{empirical
validation}.  The distinction is critical:
\begin{itemize}
  \item \textbf{What our simulations prove}:
    The axioms (C1, C2, C3) are mutually consistent and the
    derived theorems (Television Theorem, Soul Decomposition,
    Two-Kill Mechanism) follow logically.
  \item \textbf{What our simulations do NOT prove}:
    That real brains obey our axioms.
    That $\Gamma_{\text{meta}}$ corresponds to subjective
    experience.  That $K\approx 128$ for cortical pyramidal
    cells.  That the thalamic gate equation is quantitatively
    correct.
\end{itemize}

In the language of Popper: our theory is \emph{falsifiable}
(see below) but has not yet been \emph{tested}.  The zero
failure rate across all simulations is not evidence of
correctness---it is evidence of circularity.
Only external experiments on biological tissue, conducted by
independent groups with no stake in the theory, can
promote these results from mathematical consistency to
physical truth.

\subsection{Falsifiable predictions}

\begin{enumerate}
  \item \textbf{K-dependent thermal vulnerability}:
    Brain regions with higher axon dimension~$K$ should show
    earlier functional degradation during controlled hyperthermia.
    PFC ($K\sim 128$) should degrade before V1 ($K\sim 8$).
  \item \textbf{Arousal--gate multiplicativity}:
    Partial brainstem lesions reducing $\Zarousal$ to 50\%
    should reduce \emph{all} modality gate gains equally
    (multiplicative, not additive).
  \item \textbf{Startle latency ordering}:
    During rapid loss of consciousness (impact, syncope),
    auditory and vestibular percepts (buzzing, spinning)
    should be reported as the \emph{last} subjective
    experiences before blackout, because the startle
    circuit is a hardware bypass.
  \item \textbf{Feedback disruption test}:
    Selectively disrupting efference copy pathways
    (e.g., via TMS to parietal cortex) should reduce
    $\Phimeta$ without reducing $\Zarousal$, $G_i$,
    or $\QPFC$---a third kill mechanism that disables the
    feedback loop itself.
\end{enumerate}

\subsection{Relation to prior work}

The consciousness index $\Phimeta$ is analogous to
Tononi's Integrated Information Theory ($\Phi$) but differs
in two critical ways:
(i)~$\Phimeta$ is defined operationally via impedance
reflection, not information-theoretic integration;
(ii)~$\Phimeta$ has an explicit necessary condition
(the television inequality) that IIT lacks.
The soul decomposition $\mathcal{S}=\Acut+Z_0^{\text{DNA}}+\rho_{\text{irred}}^2$
is, to our knowledge, the first formal definition that is
both physically grounded and empirically invariant.


\subsection{Infantile amnesia: the SNR floor of episodic memory}
\label{sec:infantile-amnesia}

Humans reliably fail to recall autobiographical events
from before age~$\sim$3 (``infantile amnesia'').
Classical explanations invoke hippocampal immaturity or
language development.  The impedance framework provides a
quantitative alternative.

The signal-to-noise ratio for memory \emph{retrieval}
is:
\begin{equation}
  \mathrm{SNR}_{\text{mem}} =
    \frac{|\Delta Z_{\text{mem}}|^2}
         {\sigma^2_{Z_{\text{bg}}}}\,,
  \label{eq:SNR-mem}
\end{equation}
where
$|\Delta Z_{\text{mem}}|^2$ is the impedance change
written by a single episode (numerator) and
$\sigma^2_{Z_{\text{bg}}}$ is the background variance of
the impedance map (denominator).

Before age~3, three processes simultaneously inflate
$\sigma^2_{Z_{\text{bg}}}$:
\begin{enumerate}
  \item \textbf{Myelination}: ongoing myelination
    shifts the characteristic impedance
    $Z_0^{\text{axon}}$ on a daily basis,
    moving the reference plane against which all
    $\Gamma$ values are defined.
    A memory encoded against yesterday's $Z_0$ is
    unrecoverable when $Z_0$ has drifted.
  \item \textbf{Synaptic pruning}: the neonatal synapse
    count is approximately twice the adult value;
    pruning eliminates $\sim$50\% of synapses,
    deleting the $Z_k$ channels on which early
    $\Delta Z$ were stored.
  \item \textbf{PFC immaturity}: with $Q_{\text{PFC}}\ll 1$
    (Sec.~\ref{sec:consciousness}), the self-referential
    loop quality is too low to generate the meta-tag
    ``\emph{I} experienced this''---the necessary encoding
    for an \emph{episodic} (rather than procedural) memory.
\end{enumerate}

Crucially, the infant \emph{does} form Hebbian traces
(C2 operates from birth with maximal $\eta$).
Motor memories ($\Gamma\to 0$; walking, grasping) survive
because they are encoded as impedance \emph{structure},
not as retrievable \emph{reflections}.
It is only declarative/episodic memory
($\Gamma\approx\varepsilon>0$, Sec.~\ref{sec:memory})
that requires $\mathrm{SNR}_{\text{mem}}>\theta_{\text{recall}}$.

\begin{corollary}[Memory requires a stable reference plane]
Episodic memory formation requires that the background
impedance variance $\sigma^2_{Z_{\text{bg}}}$ be small
relative to the single-episode encoding strength
$|\Delta Z_{\text{mem}}|^2$.
Since $\sigma^2_{Z_{\text{bg}}}$ decreases monotonically
with age during development (as myelination completes and
pruning subsides), the onset of stable episodic memory
at age~$\sim$3 corresponds to the epoch when
$\sigma^2_{Z_{\text{bg}}}$ falls below the retrieval
threshold.
\end{corollary}

\subsection{Nociception without memory: the three-layer
  pain architecture}
\label{sec:infant-pain}

A critical clarification emerges: newborn infants
\emph{do} exhibit pain responses (flexion withdrawal,
crying, cortisol elevation) despite having
$\mathrm{SNR}_{\text{mem}}\approx 0$.
This is not a contradiction.

Pain processing involves three physically distinct
circuits ordered by their SNR requirements:
\begin{enumerate}
  \item \textbf{Nociception} (spinal reflex arc):
    a hardware detector of local $|\Gamma_{\text{tissue}}|>\theta_{\text{noci}}$.
    Requires no calibrated VNA, no PFC, no consciousness.
    Fully operational at birth.
  \item \textbf{Pain perception} (thalamo-cortical loop):
    requires $\Phi_{\text{meta}}>0$, i.e., the television
    inequality must be satisfied.
    Partially operational in infancy (low $Q_{\text{PFC}}$).
  \item \textbf{Pain memory} (hippocampal-PFC encoding):
    requires $\mathrm{SNR}_{\text{mem}}>\theta_{\text{recall}}$.
    Non-operational before age~$\sim$3.
\end{enumerate}

The neonatal design thus prioritises \emph{survival}
(hardware interrupt) over \emph{experience}
(consciousness) over \emph{record-keeping} (memory),
bringing each layer online only when the preceding
layer's reference plane has stabilised.

\subsection{Mechanical percussion and consciousness
  restoration}
\label{sec:percussion}

A universally recognised but rarely analysed phenomenon:
when a CRT television displays a degraded signal, a
physical slap to the chassis often temporarily restores
picture quality.
Humans exhibit an identical behaviour---slapping one's
own forehead to ``clear the head'' during drowsiness or
cognitive fog.

The physics is the same in both systems.
In an aging television, oxidised solder joints and
cold-solder cracks create intermittent contact impedance:
$\Gamma_{\text{contact}}>0$, fluctuating near the
conduction threshold.
Mechanical vibration transiently re-seats the contact
surfaces, reducing $\Gamma_{\text{contact}}\to 0$ for
seconds to minutes.

In the brain, the analogous ``cold-solder cracks'' are:
\begin{itemize}
  \item \textbf{Adenosine-mediated synaptic fatigue}:
    adenosine accumulation shifts synaptic
    $Z_{\text{synapse}}$ away from its matched
    operating point ($\Gamma_{\text{syn}}\uparrow$).
  \item \textbf{Thalamic gate drift}: with fatigue or
    hypoglycemia, gate gains $G_i$ drift below threshold.
  \item \textbf{Arousal fluctuation}: sub-threshold
    $G_{\text{arousal}}$ causes
    $\Phi(t)=G_{\text{arousal}}\cdot\sum G_i\cdot Q_{\text{PFC}}$
    to hover near~$\theta$, producing intermittent
    consciousness (microsleeps, ``spacing out'').
\end{itemize}

A physical slap to the head produces three simultaneous
effects that transiently restore $\Phi>\theta$:
\begin{enumerate}
  \item \textbf{Cerebrospinal fluid pressure wave}:
    mechanical perturbation reaches synaptic junctions,
    transiently resetting $\Gamma_{\text{syn}}$ closer
    to the matched position---identical to the television
    mechanism.
  \item \textbf{Nociceptive arousal kick}:
    the slap is a mild nociceptive stimulus;
    trigeminal afferents activate the brainstem
    reticular activating system, directly elevating
    $G_{\text{arousal}}$.
  \item \textbf{Vestibular startle bypass}:
    the head acceleration triggers the vestibular
    startle circuit (Sec.~\ref{sec:two-kills}),
    which forces $G_{\text{vestibular}}=1$ via
    hardware interrupt, injecting a strong signal
    into $\sum G_i$.
\end{enumerate}

The combined effect pushes $\Phi$ above $\theta$ for
a transient period---exactly as a slap to the television
restores picture quality until the contact re-oxidises.
The critical difference: the brain has $\Gamma_{\text{meta}}$
(closed-loop self-reference), so the individual is
\emph{aware} that clarity has been restored;
the television is not.

\begin{remark}[Not metaphor but identity]
The television analogy (Theorem~\ref{thm:TV-full}) is not
a didactic convenience.
The percussion-restoration effect in both systems arises
from the same impedance physics: mechanical perturbation
transiently reduces contact $\Gamma^2$ at degraded
junctions.
The identity of the repair mechanism across biological
and electronic systems is a direct consequence of the
universality of $\Gamma$ (Theorem~\ref{thm:universality}
in Paper~VI).
\end{remark}

\subsection{The $\Gamma$-Pain Index: objective nociceptive
  burden without consciousness}
\label{sec:gamma-pain-index}

The three-layer architecture (Sec.~\ref{sec:infant-pain})
exposes a critical gap in current clinical practice.
The standard-of-care pain assessment---a subjective
numeric rating scale (NRS, 0--10)---measures only the
\emph{Layer~3 verbal report}, which is a nonlinear,
non-monotonic, and individually variable function of the
underlying tissue state:
\begin{equation}
  \text{Pain\_Score} =
    g\!\left(
      \Phi_{\text{meta}}\cdot|\Gamma_{\text{tissue}}|
      + \varepsilon_{\text{noise}}
    \right),
  \label{eq:pain-score-function}
\end{equation}
where $g(\cdot)$ encodes psychological defence mechanisms,
social context, language ability, and arousal state.
Critically:
\begin{itemize}
  \item $\text{Pain\_Score}=0$ does not imply
    $\Gamma_{\text{tissue}}=0$
    (the patient may be sedated, aphasic, or stoic);
  \item $\text{Pain\_Score}=10$ does not imply
    $|\Gamma|\to 1$
    (anxiety, catastrophising, or social amplification
    can inflate $g$ at moderate $\Gamma$).
\end{itemize}

For populations where Layer~3 is unreliable or
non-existent---neonates, sedated ICU patients, persons
with advanced dementia, aphasic stroke survivors---the
NRS provides \emph{zero information}.
Yet the Layer~1 nociceptive signal
($|\Gamma_{\text{tissue}}|>\theta_{\text{noci}}$)
continues to operate, driving autonomic responses,
stress hormone release, and cumulative impedance debt
($D_Z$; Paper~III).

We define the \textbf{$\Gamma$-Pain Index}
$\Pi_{\Gamma}(t)$ as a consciousness-independent,
objective measure of nociceptive burden:
\begin{equation}
  \boxed{\;
  \Pi_{\Gamma}(t) =
    w_n\,\Gamma_n^2(t) + w_v\,\Gamma_v^2(t)
    + w_D\,\frac{dD_Z}{dt}\;
  }
  \label{eq:gamma-pain-index}
\end{equation}
where:
\begin{itemize}
  \item $\Gamma_n^2(t)$: neural-side mismatch, estimated
    from EEG spectral changes (delta/theta power shift,
    pain-related evoked potential amplitude, coherence
    disruption between somatosensory and frontal regions);
  \item $\Gamma_v^2(t)$: vascular-side mismatch, estimated
    from HRV (LF/HF ratio departure from baseline),
    blood pressure variability, skin conductance, and
    regional blood flow;
  \item $dD_Z/dt$: instantaneous rate of impedance debt
    accumulation (Paper~III), capturing the
    Coffin--Manson fatigue cost of the current stressor;
  \item $w_n, w_v, w_D$: organ-specific weighting
    coefficients (derivable from the dual-network model
    of Paper~II).
\end{itemize}

The critical property of $\Pi_{\Gamma}$ is that
\emph{it does not depend on $\Phi_{\text{meta}}$}.
Consciousness modulates the \emph{experience} of pain
(Layer~2) and the \emph{report} of pain (Layer~3), but
it does not modulate the tissue-level impedance mismatch
(Layer~1) that $\Pi_{\Gamma}$ measures.
A patient under general anaesthesia with
$\Phi_{\text{meta}}\approx 0$ and
$\text{Pain\_Score}=\text{undefined}$ can nonetheless
have $\Pi_{\Gamma}\gg 0$, indicating significant
nociceptive load and ongoing impedance debt accumulation.

\begin{table}[h]
\caption{Comparison of pain assessment modalities by layer.}
\label{tab:pain-layers}
\begin{ruledtabular}
\begin{tabular}{lcccc}
\textbf{Modality} & \textbf{Layer} &
  \textbf{Requires} & \textbf{Fails in} &
  \textbf{Measures} \\
  & & \textbf{consciousness?} & & \\
\hline
NRS (0--10) & L3 & Yes & neonate, sedated, & subjective \\
  & & & aphasic, demented & report \\
FLACC/BPS & L2--L3 & Partial & deeply sedated, & behavioural \\
  & & & paralysed & proxy \\
HR/HRV/BP & L1--L2 & No & $\beta$-blocked, & autonomic \\
  & & & paced & proxy \\
$\Pi_{\Gamma}$ & L1 & No & none$^*$ & tissue \\
  & & & & mismatch \\
\end{tabular}
\end{ruledtabular}
{\footnotesize $^*$Requires multi-modal sensor input
(EEG + HRV + bioimpedance); no dependence on patient
cooperation or consciousness.}
\end{table}

\begin{remark}[Falsifiable prediction]
In sedated ICU patients receiving identical surgical
stimuli, $\Pi_{\Gamma}$ (computed from EEG + HRV)
should correlate with tissue damage severity
(wound area, inflammatory marker levels) with
$r>0.5$, \emph{independent} of anaesthetic depth
(BIS score).  A null result---no correlation between
$\Pi_{\Gamma}$ and tissue damage when consciousness
is controlled---would falsify the Layer~1/Layer~3
separation.
\end{remark}

\begin{remark}[Clinical protocol: multi-layer pain assessment]
The $\Gamma$-Pain Index motivates a four-tier clinical
protocol:
\begin{enumerate}
  \item \textbf{Subjective report} (NRS/VAS):
    when available, respect the patient's Layer~3 output.
  \item \textbf{Behavioural observation}
    (FLACC/BPS/CPOT): capture Layer~2 when Layer~3 is
    unavailable.
  \item \textbf{Autonomic proxies}
    (HR, HRV, BP, skin conductance): Layer~1--2 surrogates
    available in any monitored patient.
  \item \textbf{$\Pi_{\Gamma}$ composite}
    (EEG + HRV + bioimpedance): direct Layer~1
    measurement, the only modality that reports tissue
    mismatch \emph{regardless of consciousness state}.
\end{enumerate}
Relying on Layer~3 alone---asking ``how much does it
hurt?'' as the sole assessment---discards the majority
of the information and is physically equivalent to
diagnosing a transmission-line fault by asking the
television whether the picture looks good.
\end{remark}


% ============================================================
\section{Conclusion}
\label{sec:conclusion}

Starting from the Minimum Reflection Principle
$A[\Gamma]=\int\sum\Gamma^2\,dt\to\min$ and three
constraints (C1--C3), we have derived:

\begin{enumerate}
  \item \textbf{Memory} $=\partial Z/\partial t$:
    the time derivative of impedance under Hebbian update.
    Muscle memory is $\Gamma\to 0$ (frozen match);
    cognitive memory is $\Gamma\approx\varepsilon>0$
    (detectable reflection).
  \item \textbf{Consciousness}:
    $\Phimeta=1-\Gmeta^2$, the self-referential reflection
    quality of a closed-loop impedance network.
    Necessary condition:
    $\Zarousal\cdot\sum G_i\cdot\QPFC>\theta$
    (the television inequality).
    The brain and a television have the same components;
    they differ by one feedback arc.
  \item \textbf{Soul}:
    $\mathcal{S}=\ker(\partial Z/\partial t)=\Acut+Z_0^{\text{DNA}}+\rho_{\text{irred}}^2$,
    the invariant geometric structure that no experience can
    modify.
\end{enumerate}

These are not metaphors.
They are operational definitions derived from a single
variational principle with three constraints.
The physics is elementary (impedance matching, thermal noise,
linear algebra); the consequences---that memory requires
imperfection, that consciousness requires topology, and that
the soul is a null space---are not.


% ============================================================
\begin{acknowledgments}
This work used no external funding.
All experiments were conducted on the Alice Gamma-Net simulator.
The author thanks Claude (Anthropic) for extended
theoretical discussions that sharpened the Television Theorem
and the two-kill-mechanism analysis.
\end{acknowledgments}


% ============================================================
\appendix

\section{Experimental Code Availability}
\label{app:code}

The computational experiments described in
Section~\ref{sec:two-kills} are available at:
\begin{itemize}
  \item \texttt{experiments/exp\_prefrontal\_signal\_trace.py}:
    Visual and auditory signal trace to PFC.
  \item \texttt{experiments/exp\_consciousness\_two\_kills.py}:
    Fever vs.\ occipital impact comparison.
\end{itemize}
Full source code: \url{https://github.com/cyhuang76/alice-gamma-net}

DOI (papers): \href{https://doi.org/10.5281/zenodo.18751831}{10.5281/zenodo.18751831}

DOI (code): \href{https://doi.org/10.5281/zenodo.18852835}{10.5281/zenodo.18852835}


\end{document}
